\documentclass[11pt]{article}
\usepackage[margin=1in]{geometry}
\usepackage{amsmath,amssymb,amsthm,mathtools}
\usepackage{graphicx}
\usepackage[hidelinks]{hyperref}
\usepackage{xcolor}

\newtheorem{theorem}{Theorem}
\newtheorem{proposition}{Proposition}
\newcommand{\R}{\mathbb{R}}
\newcommand{\PhiL}{\Phi_{\lambda,Q}}
\newcommand{\Sph}{\mathbb{S}}

\title{A Positive-$\gamma$ Counterexample to the Nonlocal Poincar\'e Question as Stated}
\author{Automated mathematics run \texttt{fa\_banach\_001}\\
Lane 12, agent \texttt{agent\_lane\_12}, model GPT5.6}
\date{11 August 2026}

\begin{document}
\maketitle

\begin{center}
\fbox{\parbox{0.91\linewidth}{\textbf{Status.}
Candidate counterexample, likely valid.  It gives a complete negative answer
to the positive-$\gamma$ branch under the natural interpretation that the
displayed estimate is uniform for all $\lambda>0$.  The source does not print
that quantifier explicitly; a sufficiently-large-$\lambda$ repair is not
addressed.}}
\end{center}

\section{Source question and normalization}

The source is Hoai-Minh Nguyen, \emph{Characterizations of the Sobolev norms
and the total variation via nonlocal functionals, and related problems},
arXiv:2507.00842v2~\cite{Nguyen}.  Open Question~3 appears at the end of
Section~9 on PDF page~25.

\begin{center}
\includegraphics[width=\textwidth]{figures/open_problem_crop.png}
\end{center}

For a unit cube $Q\subset\R^N$, the domain-correct normalized local
functional is
\begin{equation}\label{eq:local-functional}
 \Phi_{\lambda,Q}(u)
 :=\lambda^p\!\iint_{\substack{Q\times Q\\
 |u(x)-u(y)|>\lambda |x-y|^{1+\gamma/p}}}
 |x-y|^{-N+\gamma}\,dx\,dy .
\end{equation}
The question asks, for $p\ge1$ and
$\gamma\in\R\setminus(-1,0]$ with $\kappa_{N,p,\gamma}>0$, whether
\begin{equation}\label{eq:question}
 \iint_{Q\times Q}|u(x)-u(y)|^p\,dx\,dy
 \le C_{N,p}\bigl(\Phi_{\lambda,Q}(u)+\lambda^p\bigr),
 \qquad u\in L^p(Q).
\end{equation}
The source displays a free $\lambda$ and a constant independent of $\lambda$;
we use the standard uniform-in-$\lambda>0$ reading.  Nguyen writes
$\Phi_\lambda(u)$ in the question even though $u$ is defined only on $Q$.
Equation~\eqref{eq:local-functional} is the local functional defined in the
paper.  A global definition in the TeX source accidentally omits the factor
$\lambda^p$, but the local definition, Theorem~1.3, and all later calculations
use this factor.

\section{Proof intuition}

For $\gamma>0$ the kernel $|x-y|^{-N+\gamma}$ is integrable across the
diagonal on a bounded cube.  Consequently the entire normalized nonlocal
energy is at most a finite constant times $\lambda^p$, no matter what $u$ is.
Sending $\lambda$ to zero therefore erases the proposed right-hand side, while
the mean oscillation of any nonconstant $u$ remains positive.

There is one potential logical escape: perhaps no positive $\gamma$ satisfies
the hypothesis $\kappa_{N,p,\gamma}>0$.  Gobbino and Picenni's
$\Gamma$-liminf theorem~\cite{GobbinoPicenni} closes that escape.  A clipping
argument converts Nguyen's convergence-in-measure cell problem to their
$L^1$ topology without increasing the energy, proving that the cell constant
is positive for every $\gamma>0$.

\section{The counterexample}

\begin{theorem}[Positive-$\gamma$ branch fails]\label{thm:main}
Let $N\ge1$, $p\ge1$, $\gamma>0$, and $Q=(0,1)^N$.  Then
$\kappa_{N,p,\gamma}>0$, but there is no finite constant $C$ for which
\eqref{eq:question} holds for every $u\in L^p(Q)$ and every $\lambda>0$.
Indeed the assertion fails for $u(x)=x_1$.  This remains true if $C$ is
allowed to depend on $N,p,\gamma$, and $Q$.
\end{theorem}

We separate verification of the premise from the short counterexample.

\begin{proposition}[The positive-$\gamma$ premise is nonvacuous]
For every $N\ge1$, $p\ge1$, and $\gamma>0$, Nguyen's cell constant satisfies
\begin{equation}\label{eq:kappa-lower}
 \kappa_{N,p,\gamma}
 \ge A_{N,p}\frac{\log 2}{2^{\gamma+1}-1}>0,
\end{equation}
where
$A_{N,p}=\int_{\Sph^{N-1}}|v\cdot\omega|^p\,
d\mathcal H^{N-1}(\omega)>0$ for any unit vector $v$.
\end{proposition}

\begin{proof}
In Nguyen's definition, one takes the infimum of
$\liminf_{\lambda\to\infty}\Phi_{\lambda,Q}(h_\lambda)$ over measurable
families converging in measure on $Q$ to
\[
 h(x)=\frac{x_1+\cdots+x_N}{\sqrt N}.
\]
Fix such a family and let
$T(t)=\min\{\max\{t,0\},\sqrt N\}$.  Since $T$ is $1$-Lipschitz and
$T(h)=h$, its exceedance set is contained in that of the original function:
\[
 \Phi_{\lambda,Q}(T\circ h_\lambda)
 \le \Phi_{\lambda,Q}(h_\lambda).
\]
Moreover $T\circ h_\lambda\to h$ in measure.  The functions are uniformly
bounded on the finite-measure cube, so this convergence is also in $L^1(Q)$.

Gobbino--Picenni's Theorem~1.1~\cite{GobbinoPicenni}, applied on the open cube,
now gives
\[
 \liminf_{\lambda\to\infty}
 \Phi_{\lambda,Q}(T\circ h_\lambda)
 \ge A_{N,p}\frac{\log2}{2^{\gamma+1}-1}F_p(h,Q).
\]
Here $F_p(h,Q)=1$: the gradient of $h$ has Euclidean norm one and $|Q|=1$
(the same computation is the total variation when $p=1$).  Combining the two
inequalities and taking the infimum over all families proves
\eqref{eq:kappa-lower}.
\end{proof}

\begin{proof}[Proof of Theorem~\ref{thm:main}]
Set $u(x)=x_1$.  By Fubini,
\begin{align}\label{eq:lhs}
 \iint_{Q\times Q}|u(x)-u(y)|^p\,dx\,dy
 &=\int_0^1\!\int_0^1|s-t|^p\,ds\,dt\\
 &=2\int_0^1\!\int_0^s(s-t)^p\,dt\,ds
 =\frac{2}{(p+1)(p+2)}>0.\notag
\end{align}
On the other hand, because $Q-Q\subset B_{\sqrt N}(0)$ and $\gamma>0$,
\begin{align}\label{eq:kernel}
 \nu_\gamma(Q\times Q)
 &=\iint_{Q\times Q}|x-y|^{-N+\gamma}\,dx\,dy\\
 &\le |\Sph^{N-1}|\int_0^{\sqrt N}r^{\gamma-1}\,dr
 =\frac{|\Sph^{N-1}|}{\gamma}N^{\gamma/2}<\infty.\notag
\end{align}
The exceedance set in~\eqref{eq:local-functional} is contained in $Q\times Q$,
so
\begin{equation}\label{eq:vanish}
 0\le\Phi_{\lambda,Q}(u)
 \le \frac{|\Sph^{N-1}|}{\gamma}N^{\gamma/2}\lambda^p.
\end{equation}
Thus $C(\Phi_{\lambda,Q}(u)+\lambda^p)\to0$ as
$\lambda\downarrow0$ for every finite $C$, contradicting the fixed positive
quantity~\eqref{eq:lhs}.  The preceding proposition verifies the required
hypothesis $\kappa_{N,p,\gamma}>0$.
\end{proof}

\section{Global zero-extension robustness}

Although the local functional is the natural reading, the obstruction also
survives a normalized global interpretation.  Extend $u(x)=x_1$ on $Q$ by
zero and call the result $U$ on $\R^N$.  If $(x,y)$ belongs to its exceedance
set, at least one endpoint lies in $Q$, $|U(x)-U(y)|\le1$, and therefore
\[
 |x-y|<R_\lambda:=\lambda^{-p/(p+\gamma)}.
\]
Double-counting which endpoint lies in $Q$ gives
\begin{align*}
 \Phi_\lambda(U)
 &\le2\lambda^p\int_Q\int_{|x-y|<R_\lambda}
 |x-y|^{-N+\gamma}\,dy\,dx\\
 &=\frac{2|\Sph^{N-1}|}{\gamma}
 \lambda^pR_\lambda^\gamma
 =\frac{2|\Sph^{N-1}|}{\gamma}
 \lambda^{p^2/(p+\gamma)}\longrightarrow0.
\end{align*}
Hence zero extension does not rescue the estimate.

\section{Verification, scope, and novelty}

The argument is analytic and exact.  It checks the premise separately, gives
an explicit positive left-hand side, and bounds the entire energy rather than
approximating its exceedance set.  It also proves the stronger statement that
no additive remainder $r(\lambda)\to0$ can make a uniform inequality of this
form true for positive $\gamma$.

The conclusion is deliberately scoped.  It answers the positive-$\gamma$
branch of Open Question~3 under the uniform reading forced by the free
parameter in the display.  It does not address $\gamma\le-1$.  If the intended
question was instead restricted to all sufficiently large $\lambda$ when
$\gamma>0$, the present small-$\lambda$ argument does not apply.

A bounded novelty search on 11 August 2026 covered the run indexes, the exact
arXiv id and title, exact-phrase searches for Open Question~3, searches for
``nonlocal Poincar\'e'', $\Phi_\lambda$, and $\kappa_{N,p,\gamma}$, and the
source paper's references.  No later paper explicitly stating this
counterexample or resolving Open Question~3 was found.  This supports but
does not certify novelty.  Human review should focus on the intended
$\lambda$ quantifier and the source's local/global notation.

\begin{thebibliography}{9}
\bibitem{Nguyen}
H.-M. Nguyen,
\emph{Characterizations of the Sobolev norms and the total variation via
nonlocal functionals, and related problems},
arXiv:2507.00842v2, 2025.

\bibitem{GobbinoPicenni}
M. Gobbino and N. Picenni,
\emph{Gamma-liminf estimate for a class of non-local approximations of Sobolev
and BV norms},
J. Funct. Anal. \textbf{289} (2025), no.~9, Paper 111106;
arXiv:2311.05560.
\end{thebibliography}

\end{document}

